%% LyX 2.4.4 created this file.  For more info, see https://www.lyx.org/.
%% Do not edit unless you really know what you are doing.
\documentclass[english]{article}
\usepackage[T1]{fontenc}
\usepackage[latin9]{inputenc}
\usepackage{units}
\usepackage{enumitem}
\usepackage{amsmath}
\usepackage{amsthm}
\usepackage{graphicx}
\usepackage{geometry}
\geometry{verbose,tmargin=1in,bmargin=1in,lmargin=1in,rmargin=1in}

\makeatletter
%%%%%%%%%%%%%%%%%%%%%%%%%%%%%% Textclass specific LaTeX commands.
\numberwithin{equation}{section}
\numberwithin{figure}{section}
\newlength{\lyxlabelwidth}      % auxiliary length 

%%%%%%%%%%%%%%%%%%%%%%%%%%%%%% User specified LaTeX commands.
\usepackage{fancyhdr}
\usepackage{lastpage}
\pagestyle{fancy}
\usepackage{enumitem}
\usepackage{ifthen}
\everymath=\expandafter{\the\everymath\displaystyle}
%\DeclareMathSizes{10}{10}{10}{10}
\usepackage{multicol}

\usepackage{tikz}
\usetikzlibrary{patterns}
\usepackage{pgfplots}
% Added by lyx2lyx
\setlength{\parskip}{\smallskipamount}
\setlength{\parindent}{0pt}

\makeatother

\usepackage{babel}
\begin{document}
\renewcommand{\author}{Dr.\,James Ashe}

\newcommand{\classNum}{232}

\newcommand{\classSec}{A and B}

\newcommand{\examType}{Quiz}

\renewcommand{\date}{February 11, 2013} 

%%First page's header%%%%%%%%%%%%%%%%%%%%%%
\fancypagestyle{firststyle} %%header settings for first pagr
{
	\fancyhead[L]{MTH \classNum{}(\classSec{}) \author{}\\
	\textbf{\examType{}} \date{}}
	\fancyhead[C]{\textbf{Name:}}
	\setlength{\headsep}{14pt}
}
%%%%%%%%%%%%%%%%%%%%%%%%%%%%%%%%%%%%%%%%%%

%%Next page's header%%%%%%%%%%%%%%%%%%%%%%%
\setlist{nolistsep}
\lhead{\classNum{}(\classSec{})} 
\chead{\textbf{Name:}} 
\cfoot{\thepage{}~of~\pageref{LastPage}} 
\setlength{\headsep}{5pt} 
%%%%%%%%%%%%%%%%%%%%%%%%%%%%%%%%%%%%%%%%%%%

%%Various settings; see the preamble for other settings%%%%%
%%\setenumerate[0]{leftmargin=12pt,itemindent=0pt,labelwidth=0pt}
\renewcommand{\labelenumi}{\textbf{\arabic{enumi}.}}
\renewcommand{\labelenumii}{\textbf{\alph{enumii}.}}
\thispagestyle{firststyle}
%%%%%%%%%%%%%%%%%%%%%%%%%%%%%%%%%%%%%%%%%%

\newcommand{\xmax}{0}
\newcommand{\xmin}{-\xmax}
\newcommand{\xstep}{0}
\newcommand{\ymax}{0}
\newcommand{\ymin}{-\ymax}
\newcommand{\ystep}{0} 

\textbf{Justify your answers and show your work. }This quiz will be
worth two homework grades. 

\emph{Complete the following steps for $f(x)=24-2x$, interval $[0,6]$,
and $n=3$.}
\begin{enumerate}[resume]
\item Sketch the graph of the function on the given interval and illustrate
the \emph{right Riemann sums }using $n$ subintervals. Clearly label
the gridpoints.\\
\vspace{2.0in}
\item Calculate the \emph{right Riemann sums }for the given interval and
number of subintervals, $n$.\vspace{2cm}
\end{enumerate}
\emph{The figure shows the area of regions bounded by the graph of
$f(x)$ and the $x-$axis. Evaluate the following integrals.}

\begin{multicols}{2}
\begin{enumerate}[resume]
\item ${\displaystyle \int_{0}^{8}f}(x)\,dx$\vfill
\item ${\displaystyle \int}_{5}^{8}2f(x)\,dx$\vfill
\item ${\displaystyle \int_{0}^{5}f(x)\,dx}-{\displaystyle \int_{8}^{0}f(x)\,dx}$\vfill\columnbreak
\end{enumerate}
\includegraphics[width=3in]{D:/home/jim/JamesAshe/JCSU/Calc_2_MTH_232/images/Exam_1_area_under_graph.svg}

\end{multicols}
\begin{enumerate}[resume]
\item [\textbf{6.}]\setcounter{enumi}{6}Write the series in summation notation:
$2+4+6+8+10$.
\end{enumerate}
\newpage

\emph{Use area to evaluate the definite integral.}
\begin{enumerate}[resume]
\item ${\displaystyle \int_{-9}^{9}\sqrt{81-x^{2}}}\,dx$\vfill
\end{enumerate}
\emph{Use the value of the first integral $I$ to evaluate the given
integral. $I={\displaystyle \int_{0}^{1}\left(x^{3}-5x\right)\,dx}=-\frac{7}{4}$.}
\begin{enumerate}[resume]
\item ${\displaystyle \int_{0}^{1}\left(10x-2x^{3}\right)\,dx}$\vfill
\end{enumerate}
\emph{Evaluate the following.}
\begin{enumerate}[resume]
\item Let $a$ and $b$ be constants ${\displaystyle {\displaystyle \frac{d}{dx}\int_{a}^{x}f(t)\,dt}}$\vfill
\item $\frac{d}{dx}{\displaystyle \int_{x}^{2}\sqrt{t^{5}}}\,dt$\vfill
\item $\int_{1}^{2}(x^{2}-8x+1)\,dx$\vfill\vspace{1cm}\newpage
\item ${\displaystyle \int\frac{x-\sqrt{x}}{x^{3}}}\,dx$\vfill
\item $f(x)=\cos3x$

\begin{enumerate}
\item ${\displaystyle \int\cos3x\,dx}$\vfill\vspace{-2.0cm}
\item ${\displaystyle \int_{-\nicefrac{\pi}{3}}^{\nicefrac{\pi}{3}}\cos3x\,dx}$\vfill\vspace{-2.0cm}
\end{enumerate}
\item Find the solution to the following intital value problem. $g'(x)=4x\left(x^{3}-\frac{1}{4}\right)$
and $g(1)=2.$\vfill
\item A bowling ball rolls down a ramp at a speed of $v(t)=-3t^{2}$ feet
per second (where $t$ is seconds). The distance of the ball from
the top of the ramp is $4$ feet after $2$ seconds. Find the function
giving the distance of the ball from the top of the ramp at $t$ seconds.\vfill
\end{enumerate}

\end{document}
